% GENERATED FILE. Edit canonical lessons, then run npm run generate:books.
% canonical-source-hash: fnv1a64:fc3f4bce7e68ee23
% canonical-lessons: ES-C732-saludable, ES-C732-cuidarse, ES-C732-enfermedad, ES-C732-enfermarse, ES-C732-empeorar

\chapter{Cuidarse}
\label{ch:es-cuidarse}
\hlchaptermodality{732}

\begin{chapteropening}
\textbf{By the end of this chapter:} \emph{I can say how I look after my health and when someone gets worse.}

Complete the last lesson of chapter 732: I can say how I look after my health and when someone gets worse.
\end{chapteropening}

\section[saludable]{saludable — healthy}

\label{lesson:ES-C732-saludable}



\begin{quote}
Before the new one: say the Spanish for optimistic, then the Spanish for pessimistic.
\end{quote}

\subsection*{You'll want to know: saludable}
\textbf{saludable} — \textquotedblleft{}healthy\textquotedblright{}.

\begin{quote}
Como de forma saludable.
\end{quote}

\begin{quote}
\emph{I eat in a healthy way.}
\end{quote}

For food and habits; a healthy person is \emph{sano}.

\begin{cousinweb}
\emph{saludable} is built on \emph{salud} \textquotedblleft{}health\textquotedblright{} with the ending \emph{-able}; \emph{salud} comes from Latin \emph{salūs}, \emph{salūtem} \textquotedblleft{}health, safety\textquotedblright{} (EtymDB).
\end{cousinweb}

\begin{grammarlens}[title={What you've built}]
One more word for telling, explaining and describing at B1.
\end{grammarlens}

\subsection*{Your turn}
\begin{itemize}
\raggedright
  \item \emph{Say it:} \emph{saludable}
  \item \emph{Say it:} \emph{saludable}, once more
  \item \emph{Say it:} a sentence of your own with \emph{saludable}
  \item \emph{Recall:} say the Spanish for optimistic, then the Spanish for pessimistic, then say \emph{saludable} again
\end{itemize}

\begin{tcolorbox}[breakable,colback=teal!4,colframe=teal!35!black,title={Before you move on}]
What does saludable mean? (\textquotedblleft{}Healthy\textquotedblright{}.) Say it once more.
\end{tcolorbox}
\section[cuidarse]{cuidarse — to look after oneself}

\label{lesson:ES-C732-cuidarse}



\begin{quote}
Before the new one: say the Spanish for pessimistic, then the Spanish for healthy.
\end{quote}

\subsection*{You'll want to know: cuidarse}
\textbf{cuidarse} — \textquotedblleft{}to look after oneself\textquotedblright{}.

\begin{quote}
Cuídate mucho.
\end{quote}

\begin{quote}
\emph{Take good care of yourself.}
\end{quote}

\emph{¡Cuídate!} is a common goodbye.

\begin{cousinweb}
\emph{cuidarse} is the reflexive of \emph{cuidar}, which is inherited from Latin \emph{cōgitāre} \textquotedblleft{}to think\textquotedblright{} (EtymDB); in old Spanish \emph{cuidar} meant to think, then to watch over.
\end{cousinweb}

\begin{grammarlens}[title={What you've built}]
One more word for telling, explaining and describing at B1.
\end{grammarlens}

\subsection*{Your turn}
\begin{itemize}
\raggedright
  \item \emph{Say it:} \emph{cuidarse}
  \item \emph{Say it:} \emph{cuidarse}, once more
  \item \emph{Say it:} a sentence of your own with \emph{cuidarse}
  \item \emph{Recall:} say the Spanish for pessimistic, then the Spanish for healthy, then say \emph{cuidarse} again
\end{itemize}

\begin{tcolorbox}[breakable,colback=teal!4,colframe=teal!35!black,title={Before you move on}]
What does cuidarse mean? (\textquotedblleft{}To look after oneself\textquotedblright{}.) Say it once more.
\end{tcolorbox}
\section[la enfermedad]{la enfermedad — illness, disease}

\label{lesson:ES-C732-enfermedad}



\begin{quote}
Before the new one: say the Spanish for healthy, then the Spanish for to look after oneself.
\end{quote}

\subsection*{You'll want to know: la enfermedad}
\textbf{la enfermedad} — \textquotedblleft{}illness, disease\textquotedblright{}.

\begin{quote}
Es una enfermedad muy rara.
\end{quote}

\begin{quote}
\emph{It is a very rare disease.}
\end{quote}

\begin{cousinweb}
\emph{enfermedad} is inherited from Latin \emph{īnfirmitās} \textquotedblleft{}weakness\textquotedblright{} (EtymDB), built on \emph{īnfirmus} \textquotedblleft{}not firm, weak\textquotedblright{}.
\end{cousinweb}

\begin{grammarlens}[title={What you've built}]
One more word for telling, explaining and describing at B1.
\end{grammarlens}

\subsection*{Your turn}
\begin{itemize}
\raggedright
  \item \emph{Say it:} \emph{la enfermedad}
  \item \emph{Say it:} \emph{la enfermedad}, once more
  \item \emph{Say it:} a sentence of your own with \emph{la enfermedad}
  \item \emph{Recall:} say the Spanish for healthy, then the Spanish for to look after oneself, then say \emph{la enfermedad} again
\end{itemize}

\begin{tcolorbox}[breakable,colback=teal!4,colframe=teal!35!black,title={Before you move on}]
What does la enfermedad mean? (\textquotedblleft{}Illness, disease\textquotedblright{}.) Say it once more.
\end{tcolorbox}
\section[enfermarse]{enfermarse — to fall ill, to get sick}

\label{lesson:ES-C732-enfermarse}



\begin{quote}
Before the new one: say the Spanish for to look after oneself, then the Spanish for illness.
\end{quote}

\subsection*{You'll want to know: enfermarse}
\textbf{enfermarse} — \textquotedblleft{}to fall ill, to get sick\textquotedblright{}.

\begin{quote}
Me enfermé en las vacaciones.
\end{quote}

\begin{quote}
\emph{I got sick on holiday.}
\end{quote}

Common in Latin America; in Spain people say \emph{ponerse enfermo} or \emph{enfermar}.

\begin{cousinweb}
\emph{enfermarse} is built on \emph{enfermo} \textquotedblleft{}ill\textquotedblright{}, which is inherited from Latin \emph{īnfirmus} \textquotedblleft{}weak\textquotedblright{} (EtymDB).
\end{cousinweb}

\begin{grammarlens}[title={What you've built}]
One more word for telling, explaining and describing at B1.
\end{grammarlens}

\subsection*{Your turn}
\begin{itemize}
\raggedright
  \item \emph{Say it:} \emph{enfermarse}
  \item \emph{Say it:} \emph{enfermarse}, once more
  \item \emph{Say it:} a sentence of your own with \emph{enfermarse}
  \item \emph{Recall:} say the Spanish for to look after oneself, then the Spanish for illness, then say \emph{enfermarse} again
\end{itemize}

\begin{tcolorbox}[breakable,colback=teal!4,colframe=teal!35!black,title={Before you move on}]
What does enfermarse mean? (\textquotedblleft{}To fall ill, to get sick\textquotedblright{}.) Say it once more.
\end{tcolorbox}
\section[empeorar]{empeorar — to get worse, to worsen}

\label{lesson:ES-C732-empeorar}



\begin{quote}
Before the new one: say the Spanish for illness, then the Spanish for to fall ill.
\end{quote}

\subsection*{You'll want to know: empeorar}
\textbf{empeorar} — \textquotedblleft{}to get worse, to worsen\textquotedblright{}.

\begin{quote}
El tiempo va a empeorar.
\end{quote}

\begin{quote}
\emph{The weather is going to get worse.}
\end{quote}

The opposite of \emph{mejorar}.

\begin{cousinweb}
\emph{empeorar} is built on \emph{peor} \textquotedblleft{}worse\textquotedblright{}, which is inherited from Latin \emph{peior} (EtymDB).
\end{cousinweb}

\begin{grammarlens}[title={What you've built}]
One more word for telling, explaining and describing at B1.
\end{grammarlens}

\subsection*{Your turn}
\begin{itemize}
\raggedright
  \item \emph{Say it:} \emph{empeorar}
  \item \emph{Say it:} \emph{empeorar}, once more
  \item \emph{Say it:} a sentence of your own with \emph{empeorar}
  \item \emph{Recall:} say the Spanish for illness, then the Spanish for to fall ill, then say \emph{empeorar} again
\end{itemize}

\begin{tcolorbox}[breakable,colback=teal!4,colframe=teal!35!black,title={Before you move on}]
What does empeorar mean? (\textquotedblleft{}To get worse, to worsen\textquotedblright{}.) Say it once more.
\end{tcolorbox}
